% Part: history
% Chapter: biographies
% Section: emmy-noether

\documentclass[../../../include/open-logic-section]{subfiles}

\begin{document}

\olfileid{his}{bio}{noe}

\olsection{امی نوتر}

\olphoto{noether-emmy}{Emmy Noether}

امی نوتر (\textsc{ner}-ter) در 23 مارس 1882 در ارلانگن آلمان و در
خانواده‌ای دانشگاهی از طبقهٔ متوسطِ رو به بالا به دنیا آمد. نوتر که او
را «مادر جبر نوین» خوانده‌اند، با وجود موانع جدی در برابر آموزش زنان،
سهمی بنیادین در ریاضیات و فیزیک داشت. در آلمان آن زمان، انتظار می‌رفت
دختران جوان در هنرها آموزش ببینند و اجازهٔ ورود به مدارس مقدماتی
دانشگاه را نداشتند. بااین‌حال، نوتر پس از حضور آزاد در کلاس‌های
دانشگاه‌های گوتینگن و ارلانگن، جایی که پدرش استاد ریاضی بود، سرانجام
توانست در سال 1904، هنگامی که سیاست دانشگاه ارلانگن برای پذیرش
دانشجویان زن تغییر کرد، به‌عنوان دانشجو نام‌نویسی کند. او در سال 1907
دکترای ریاضی گرفت.

نوتر با وجود شایستگی‌هایش در طول حرفهٔ خود با مقاومت فراوانی از سوی
نهادها و همکاران روبه‌رو شد. از 1908 تا 1915 بدون دستمزد در ارلانگن
تدریس کرد. در این دوره توجه داوید هیلبرت، یکی از برجسته‌ترین
ریاضی‌دانان جهان در آن زمان، را جلب کرد و هیلبرت او را به گوتینگن
دعوت کرد. اما زنان از دستیابی به مقام استادی منع می‌شدند و نوتر فقط
می‌توانست به نام هیلبرت و باز هم بدون دستمزد درس بدهد. در همین دوره
چیزی را اثبات کرد که اکنون قضیهٔ نوتر نام دارد و هنوز در فیزیک نظری به
کار می‌رود. نوتر سرانجام در 1919 حق تدریس یافت. پاسخ نقل‌شدهٔ هیلبرت
به مقاومت ادامه‌دار از سوی همکاران دانشگاهی‌اش چنین بود: «آقایان،
سنای دانشکده حمام عمومی نیست.»

نوتر در اواخر دههٔ 1920 بر جبر مجرد تمرکز کرد و سهم او این حوزه را
دگرگون ساخت. او در اثبات‌هایش اغلب از شرط موسوم به زنجیرهٔ صعودی
استفاده می‌کرد؛ شرطی که می‌گوید هیچ زنجیرهٔ اکیداً صعودی نامتناهی از
مجموعه‌هایی معین وجود ندارد. برای نمونه، برخی ساختارهای جبری که اکنون
حلقه‌های نوتری نامیده می‌شوند این خاصیت را دارند که هیچ دنبالهٔ
نامتناهی از ایدئال‌های
$I_1 \subsetneq I_2 \subsetneq \dots$ وجود ندارد. این شرط را می‌توان
به هر ترتیب جزئی تعمیم داد. در جبر، حالت خاص ایدئال‌هایی را شامل می‌شود
که با رابطهٔ زیرمجموعه مرتب شده‌اند. همچنین می‌توان شرط دوگان زنجیرهٔ
نزولی را در نظر گرفت که در آن هر دنبالهٔ اکیداً \emph{نزولی} در یک
ترتیب جزئی سرانجام پایان می‌یابد. اگر ترتیبی جزئی شرط زنجیرهٔ نزولی را
برآورده کند، می‌توان در امتداد آن ترتیب به شیوه‌ای شبیه استقرا در
امتداد ترتیب $<$ بر~$\Nat$ استقرا کرد. چنین ترتیب‌هایی
\emph{خوش‌بنیاد} یا \emph{نوتری} نامیده می‌شوند و اصل اثبات متناظر
\emph{استقرای نوتری} نام دارد.

نوتر یهودی بود و هنگامی که نازی‌ها در سال 1933 به قدرت رسیدند، او را
از مقامش برکنار کردند. خوشبختانه نوتر توانست برای جایگاهی موقت در کالج
برین‌مار پنسیلوانیا به ایالات متحده مهاجرت کند. در دوران حضورش آنجا در
پرینستون نیز سخنرانی کرد، هرچند دانشگاه را برای زنان نامهمان‌نواز یافت
\citep[81]{Dick1981}. نوتر در سال 1935 برای برداشتن توموری در رحم تحت
عمل جراحی قرار گرفت. او بر اثر عفونت ناشی از جراحی درگذشت و در برین‌مار
به خاک سپرده شد.

\begin{reading}
برای زندگی‌نامه‌ای از نوتر، \citet{Dick1981} را ببینید. مؤسسهٔ
پریمیتر برای فیزیک نظری، سخنرانی‌های خود دربارهٔ زندگی و تأثیر نوتر را
در اینترنت در دسترس گذاشته است \citep{Perimeter2015}. اگر از خواندن
خسته شده‌اید، \emph{Stuff You Missed in History Class} پادکستی دربارهٔ
زندگی و تأثیر نوتر دارد \citep{Frey2015}. مجموعه‌آثار نوتر به زبان
اصلی آلمانی در دسترس است \citep{Noether1983}.
\end{reading}

\end{document}
